\section{总结与延伸}

\subsection{五大领域核心要点回顾}

\begin{enumerate}[leftmargin=1.5em]
    \item \textbf{Agentic 架构}（27\%）：基于 \texttt{stop\_reason} 的循环控制、子 Agent 上下文隔离、程序化 hook 保障安全
    \item \textbf{工具设计与 MCP}（18\%）：清晰的工具描述、合理的 \texttt{tool\_choice} 配置、每个 Agent 4--5 个工具的上限
    \item \textbf{Claude Code 配置}（20\%）：三级 CLAUDE.md 层次、glob 模式规则、plan mode vs direct execution
    \item \textbf{Prompt 工程}（20\%）：明确指令、边界案例 few-shot、\texttt{tool\_use} + JSON schema、Batches API
    \item \textbf{上下文管理}（15\%）：持久化 case facts、信息前置、可靠升级触发器、结构化错误传播
\end{enumerate}

\subsection{Anthropic 官方推荐学习资源}

\begin{enumerate}[leftmargin=1.5em]
    \item \textbf{Building with the Claude API}——API 基础与高级用法
    \item \textbf{Introduction to Model Context Protocol}——MCP 协议入门
    \item \textbf{Claude Code in Action}——Claude Code 实战演示
    \item \textbf{Claude 101}——Claude 基础知识
\end{enumerate}

\subsection{学习建议}

\begin{importantbox}{从实践项目入手}
每个领域都配有一个推荐实践项目。建议按以下顺序完成：
\begin{enumerate}[leftmargin=1.5em]
    \item 先完成领域一的多工具 Agent 项目（覆盖面最广）
    \item 然后是领域三的 Claude Code 配置项目（直接提升日常效率）
    \item 再完成领域二和四的工具描述与数据抽取项目
    \item 最后完成领域五的错误传播项目（整合前四个领域的知识）
\end{enumerate}
\end{importantbox}

正如原文作者所言：你不需要证书来构建生产级应用，你需要的是知识本身。掌握这五大领域的核心概念并通过实践项目内化，你就具备了成为 Claude 架构师的能力。

%% --- Fin del contenido --- %%

\end{document}
